% GENERATED FILE. Edit canonical lessons, then run npm run generate:books.
% canonical-source-hash: fnv1a64:0deda0f9bf74b7df
% canonical-lessons: FR-C365-pauvrete, FR-C365-inegalite, FR-C365-discrimination, FR-C365-discriminer, FR-C365-immigration, FR-R365-revision-358-365

\chapter{Les problèmes sociaux}
\label{ch:fr-les-problemes-sociaux}
\hlchaptermodality{365}

\begin{chapteropening}
\textbf{By the end of this chapter:} \emph{I can talk about social problems: poverty, inequality, discrimination and immigration.}

Complete the last lesson of chapter 365: I can talk about social problems: poverty, inequality, discrimination and immigration.
\end{chapteropening}

\section[la pauvreté]{la pauvreté — the poverty}

\label{lesson:FR-C365-pauvrete}



\begin{quote}
Before the new one: say the French for the prejudice, then the French for the stereotype.
\end{quote}

\subsection*{You'll want to know: la pauvreté}
\textbf{la pauvreté} — \textquotedblleft{}the poverty\textquotedblright{}.

\begin{quote}
Beaucoup de familles vivent dans la pauvreté.
\end{quote}

\begin{quote}
\emph{Many families live in poverty.}
\end{quote}

From \emph{pauvre}.

\begin{cousinweb}
\emph{pauvreté} is inherited through Old French \emph{poverté} from Latin \emph{paupertās} \textquotedblleft{}poverty\textquotedblright{} (EtymDB), from \emph{pauper} \textquotedblleft{}poor\textquotedblright{} (Lewis \& Short).
\end{cousinweb}

\begin{grammarlens}[title={What you've built}]
One more word for telling, explaining and describing at B1.
\end{grammarlens}

\subsection*{Your turn}
\begin{itemize}
\raggedright
  \item \emph{Say it:} \emph{la pauvreté}
  \item \emph{Say it:} \emph{la pauvreté}, once more
  \item \emph{Say it:} a sentence of your own with \emph{la pauvreté}
  \item \emph{Recall:} say the French for the prejudice, then the French for the stereotype, then say \emph{la pauvreté} again
\end{itemize}

\begin{tcolorbox}[breakable,colback=teal!4,colframe=teal!35!black,title={Before you move on}]
What does la pauvreté mean? (\textquotedblleft{}The poverty\textquotedblright{}.) Say it once more.
\end{tcolorbox}
\section[l'inégalité]{l'inégalité — the inequality}

\label{lesson:FR-C365-inegalite}



\begin{quote}
Before the new one: say the French for the stereotype, then the French for the poverty.
\end{quote}

\subsection*{You'll want to know: l'inégalité}
\textbf{l'inégalité} — \textquotedblleft{}the inequality\textquotedblright{}.

\begin{quote}
Les inégalités entre riches et pauvres augmentent.
\end{quote}

\begin{quote}
\emph{Inequalities between rich and poor are growing.}
\end{quote}

Often plural: \emph{les inégalités}. Feminine.

\begin{cousinweb}
\emph{inégalité} is built from \emph{in-} \textquotedblleft{}not\textquotedblright{} and \emph{égalité} \textquotedblleft{}equality\textquotedblright{}, which rests on \emph{égal}, from Latin \emph{aequālis} \textquotedblleft{}equal\textquotedblright{} (Lewis \& Short).
\end{cousinweb}

\begin{grammarlens}[title={What you've built}]
One more word for telling, explaining and describing at B1.
\end{grammarlens}

\subsection*{Your turn}
\begin{itemize}
\raggedright
  \item \emph{Say it:} \emph{l'inégalité}
  \item \emph{Say it:} \emph{l'inégalité}, once more
  \item \emph{Say it:} a sentence of your own with \emph{l'inégalité}
  \item \emph{Recall:} say the French for the stereotype, then the French for the poverty, then say \emph{l'inégalité} again
\end{itemize}

\begin{tcolorbox}[breakable,colback=teal!4,colframe=teal!35!black,title={Before you move on}]
What does l'inégalité mean? (\textquotedblleft{}The inequality\textquotedblright{}.) Say it once more.
\end{tcolorbox}
\section[la discrimination]{la discrimination — the discrimination}

\label{lesson:FR-C365-discrimination}



\begin{quote}
Before the new one: say the French for the poverty, then the French for the inequality.
\end{quote}

\subsection*{You'll want to know: la discrimination}
\textbf{la discrimination} — \textquotedblleft{}the discrimination\textquotedblright{}.

\begin{quote}
La loi interdit la discrimination à l'embauche.
\end{quote}

\begin{quote}
\emph{The law forbids discrimination in hiring.}
\end{quote}

\emph{Lutter contre la discrimination} — to fight discrimination.

\begin{cousinweb}
\emph{discrimination} is a learned borrowing from Latin \emph{discrīminātiō} \textquotedblleft{}a distinguishing\textquotedblright{} (EtymDB), from \emph{discrīmen} \textquotedblleft{}a dividing line\textquotedblright{}, built on \emph{cernere} \textquotedblleft{}to separate\textquotedblright{} (Lewis \& Short).
\end{cousinweb}

\begin{grammarlens}[title={What you've built}]
One more word for telling, explaining and describing at B1.
\end{grammarlens}

\subsection*{Your turn}
\begin{itemize}
\raggedright
  \item \emph{Say it:} \emph{la discrimination}
  \item \emph{Say it:} \emph{la discrimination}, once more
  \item \emph{Say it:} a sentence of your own with \emph{la discrimination}
  \item \emph{Recall:} say the French for the poverty, then the French for the inequality, then say \emph{la discrimination} again
\end{itemize}

\begin{tcolorbox}[breakable,colback=teal!4,colframe=teal!35!black,title={Before you move on}]
What does la discrimination mean? (\textquotedblleft{}The discrimination\textquotedblright{}.) Say it once more.
\end{tcolorbox}
\section[discriminer]{discriminer — to discriminate (against)}

\label{lesson:FR-C365-discriminer}



\begin{quote}
Before the new one: say the French for the inequality, then the French for the discrimination.
\end{quote}

\subsection*{You'll want to know: discriminer}
\textbf{discriminer} — \textquotedblleft{}to discriminate (against)\textquotedblright{}.

\begin{quote}
Cette entreprise discrimine les femmes.
\end{quote}

\begin{quote}
\emph{This company discriminates against women.}
\end{quote}

\emph{Être discriminé} — to be discriminated against.

\begin{cousinweb}
\emph{discriminer} is a learned borrowing from Latin \emph{discrīmināre} \textquotedblleft{}to divide, to separate\textquotedblright{} (Lewis \& Short), from \emph{discrīmen} \textquotedblleft{}a dividing line\textquotedblright{}, built on \emph{cernere} \textquotedblleft{}to separate\textquotedblright{}.
\end{cousinweb}

\begin{grammarlens}[title={What you've built}]
One more word for telling, explaining and describing at B1.
\end{grammarlens}

\subsection*{Your turn}
\begin{itemize}
\raggedright
  \item \emph{Say it:} \emph{discriminer}
  \item \emph{Say it:} \emph{discriminer}, once more
  \item \emph{Say it:} a sentence of your own with \emph{discriminer}
  \item \emph{Recall:} say the French for the inequality, then the French for the discrimination, then say \emph{discriminer} again
\end{itemize}

\begin{tcolorbox}[breakable,colback=teal!4,colframe=teal!35!black,title={Before you move on}]
What does discriminer mean? (\textquotedblleft{}To discriminate (against)\textquotedblright{}.) Say it once more.
\end{tcolorbox}
\section[l'immigration]{l'immigration — the immigration}

\label{lesson:FR-C365-immigration}



\begin{quote}
Before the new one: say the French for the discrimination, then the French for to discriminate.
\end{quote}

\subsection*{You'll want to know: l'immigration}
\textbf{l'immigration} — \textquotedblleft{}the immigration\textquotedblright{}.

\begin{quote}
L'immigration est un sujet de débat dans beaucoup de pays.
\end{quote}

\begin{quote}
\emph{Immigration is a subject of debate in many countries.}
\end{quote}

Feminine.

\begin{cousinweb}
\emph{immigration} is built on Latin \emph{immigrāre} \textquotedblleft{}to move into\textquotedblright{} and \emph{migrātiō} \textquotedblleft{}a moving from one place to another\textquotedblright{}, from \emph{migrāre} \textquotedblleft{}to move, to wander\textquotedblright{} (Lewis \& Short).
\end{cousinweb}

\begin{grammarlens}[title={What you've built}]
One more word for telling, explaining and describing at B1.
\end{grammarlens}

\subsection*{Your turn}
\begin{itemize}
\raggedright
  \item \emph{Say it:} \emph{l'immigration}
  \item \emph{Say it:} \emph{l'immigration}, once more
  \item \emph{Say it:} a sentence of your own with \emph{l'immigration}
  \item \emph{Recall:} say the French for the discrimination, then the French for to discriminate, then say \emph{l'immigration} again
\end{itemize}

\begin{tcolorbox}[breakable,colback=teal!4,colframe=teal!35!black,title={Before you move on}]
What does l'immigration mean? (\textquotedblleft{}The immigration\textquotedblright{}.) Say it once more.
\end{tcolorbox}
\section[(dialogue)]{Retour sur les chapitres 358–365}

\label{lesson:FR-R365-revision-358-365}



\begin{quote}
Say the French for the trade union and the strike.
\end{quote}

\subsection*{Your turn}
Say these aloud:

\begin{itemize}
\raggedright
  \item the trade union, the strike, the striker, the employee, the overtime
  \item the executive, the task, the deadline, the extension, the delivery
  \item the pay rise, the promotion, the qualification, the skill, the training
  \item stressed, exhausted, effective, the efficiency, the productivity
  \item the citizen, the citizenship, the nationality, the freedom, the equality
  \item the demonstration, to demonstrate, the demonstrator, the banner, to protest
  \item the justice, unfair, the injustice, the prejudice, the stereotype
  \item the poverty, the inequality, the discrimination, to discriminate, the immigration
\end{itemize}

\begin{tcolorbox}[breakable,colback=teal!4,colframe=teal!35!black,title={Before you move on}]
The trade union? (\textbf{le syndicat}.) The strike? (\textbf{la grève}.) The striker? (\textbf{le gréviste}.) The employee? (\textbf{le salarié}.) The overtime? (\textbf{les heures supplémentaires}.) The executive? (\textbf{le cadre}.) The task? (\textbf{la tâche}.) The deadline? (\textbf{le délai}.) The extension? (\textbf{la prolongation}.) The delivery? (\textbf{la livraison}.) The pay rise? (\textbf{l'augmentation}.) The promotion? (\textbf{la promotion}.) The qualification? (\textbf{la qualification}.) The skill? (\textbf{la compétence}.) The training? (\textbf{la formation}.) Stressed? (\textbf{stressé}.) Exhausted? (\textbf{épuisé}.) Effective? (\textbf{efficace}.) The efficiency? (\textbf{l'efficacité}.) The productivity? (\textbf{la productivité}.) The citizen? (\textbf{le citoyen}.) The citizenship? (\textbf{la citoyenneté}.) The nationality? (\textbf{la nationalité}.) The freedom? (\textbf{la liberté}.) The equality? (\textbf{l'égalité}.) The demonstration? (\textbf{la manifestation}.) To demonstrate? (\textbf{manifester}.) The demonstrator? (\textbf{le manifestant}.) The banner? (\textbf{la banderole}.) To protest? (\textbf{protester}.) The justice? (\textbf{la justice}.) Unfair? (\textbf{injuste}.) The injustice? (\textbf{l'injustice}.) The prejudice? (\textbf{le préjugé}.) The stereotype? (\textbf{le stéréotype}.) The poverty? (\textbf{la pauvreté}.) The inequality? (\textbf{l'inégalité}.) The discrimination? (\textbf{la discrimination}.) To discriminate? (\textbf{discriminer}.) The immigration? (\textbf{l'immigration}.)
\end{tcolorbox}
